\documentclass[aspectratio=169,10pt]{beamer}
\usetheme[numbering=fraction,progressbar=frametitle]{metropolis}
\usepackage{fontspec}
\setsansfont{Latin Modern Sans}
\usepackage{amsmath,amssymb,mathtools,graphicx,booktabs}
\usepackage{polyglossia}
\setdefaultlanguage{french}
\setbeamertemplate{navigation symbols}{}
\setbeamerfont{frametitle}{size=\large}
\setbeamerfont{normal text}{size=\normalsize}
\setlength{\parskip}{.45em}
\hypersetup{pdfauthor={},pdfsubject={Lecture slides},pdfkeywords={},pdftitle={Primitives et équations différentielles}}
\providecommand{\N}{\mathbb{N}}
\providecommand{\Z}{\mathbb{Z}}
\providecommand{\D}{\mathbb{D}}
\providecommand{\Q}{\mathbb{Q}}
\providecommand{\R}{\mathbb{R}}
\providecommand{\C}{\mathbb{C}}
\providecommand{\vect}[1]{\overrightarrow{#1}}
\providecommand{\vi}{\vec{\imath}}
\providecommand{\vj}{\vec{\jmath}}
\providecommand{\vk}{\vec{k}}
\providecommand{\norm}[1]{\left\lVert #1\right\rVert}
\providecommand{\vc}[2]{\begin{pmatrix}#1\\#2\end{pmatrix}}
\providecommand{\deter}[2]{\det\left(#1,#2\right)}
\providecommand{\abs}[1]{\left\lvert #1\right\rvert}
\providecommand{\dd}{\,\mathrm{d}}
\providecommand{\e}{\mathrm{e}}
\providecommand{\Card}{\operatorname{Card}}
\providecommand{\E}{\operatorname{E}}
\providecommand{\V}{\operatorname{V}}
\providecommand{\degre}{^{\circ}}
\providecommand{\un}[1]{\,\mathrm{#1}}
\providecommand{\sci}[2]{#1\times10^{#2}}
\providecommand{\pgcd}{\operatorname{PGCD}}
\providecommand{\ppcm}{\operatorname{PPCM}}
\title{Primitives et équations différentielles}
\subtitle{Cours et exemples}
\date{}
\author{}
\begin{document}
\begin{frame}[plain]
\vfill
{\usebeamerfont{title}\inserttitle\par}
\medskip
{\usebeamerfont{subtitle}\insertsubtitle\par}
\medskip
{\usebeamercolor[fg]{progress bar}\rule{\textwidth}{.5pt}}
\vfill
\end{frame}
\begin{frame}{Objectifs}
Calculer des primitives ; résoudre une équation différentielle avec condition initiale ; interpréter une approximation d'Euler.
\end{frame}
\begin{frame}{Primitives sur un intervalle}
On admet à ce stade que toute fonction continue sur un intervalle admet des primitives ; le lien avec l'aire sera établi au chapitre 8.
Une primitive de \(f\) sur un intervalle \(I\) est une fonction dérivable \(F\) telle que \(F'=f\) sur \(I\). Si \(F\) est une primitive, toutes les primitives sur \(I\) sont \(F+C\), \(C\in\R\).
\end{frame}
\begin{frame}{Primitives sur un intervalle}
Si \(F'=G'=f\), alors \((F-G)'=0\) sur l'intervalle ; \(F-G\) est donc constante. Réciproquement, \((F+C)'=f\). L'hypothèse d'intervalle est essentielle : sur une réunion d'intervalles disjoints, comme \(\mathbb{R}\setminus\{0\}\), réunion de deux intervalles, la constante peut différer d'un intervalle à l'autre.
\end{frame}
\begin{frame}{Primitives sur un intervalle}
\[\begin{array}{c|c}
f(x)&F(x)\\\hline
x^n\ (n\in\Z,\ n\ne-1)&\dfrac{x^{n+1}}{n+1}\\
\dfrac{1}{x}&\ln|x|\quad\text{sur un intervalle évitant }0\\
\dfrac{1}{\sqrt{x}}&2\sqrt{x}\quad(x>0)\\
e^x&e^x\\
\cos x&\sin x\\
\sin x&-\cos x
\end{array}\]
La constante additive peut toujours être ajoutée. Les deux dernières lignes seront réinvesties après l'étude de la trigonométrie.
\end{frame}
\begin{frame}{Reconnaître une composition}
Pour une fonction \(u\) dérivable, les primitives de \(u'u^n\), \(u'e^u\), \(\dfrac{u'}{u}\) sont respectivement \(\dfrac{u^{n+1}}{n+1}\) (\(n\ne-1\)), \(e^u\), \(\ln|u|\), sur les intervalles où les expressions ont un sens. La linéarité permet de traiter les sommes.
\end{frame}
\begin{frame}{Reconnaître une composition}
Une primitive de \(6x(x^2+1)^2\) est \((x^2+1)^3\) ; une primitive de \(e^{3x}\) est \(\dfrac{e^{3x}}{3}\). Pour imposer \(F(0)=2\) dans ce dernier cas, prendre \(F(x)=\dfrac{e^{3x}}{3}+\dfrac{5}{3}\).
\end{frame}
\begin{frame}{Équations différentielles}
Une équation différentielle relie une fonction inconnue à ses dérivées. Une solution doit être vérifiée sur un intervalle précisé.
\end{frame}
\begin{frame}{Équations différentielles}
\(y'=ay\) a pour solutions \(y(x)=Ce^{ax}\). Pour \(a\ne0\), \(y'=ay+b\) a pour solutions
\[y(x)=Ce^{ax}-\dfrac ba.\]
Si \(a=0\), l'équation devient \(y'=b\) et \(y(x)=bx+C\).
Si \(g\) est une solution particulière connue de \(y'=ay+f(x)\), les solutions sont \(y=g+Ce^{ax}\).
\end{frame}
\begin{frame}{Équations différentielles}
Si \(y\) est solution, \(z(x)=e^{-ax}y(x)\) vérifie \(z'=e^{-ax}(y'-ay)=0\) ; \(z=C\). Réciproquement, la dérivée de \(Ce^{ax}\) est \(aCe^{ax}\). Pour \(y'=ay+b\), soustraire la solution constante \(-\dfrac{b}{a}\) ramène à l'équation homogène. Pour un second membre variable, faire de même avec \(g\).
\end{frame}
\begin{frame}{Équations différentielles}
\(y'=-2y+6\) et \(y(0)=1\) donnent \(y=3+Ce^{-2x}\), puis \(C=-2\). La solution \(3-2e^{-2x}\) croît vers l'équilibre \(3\) lorsque \(x\to+\infty\).
\end{frame}
\begin{frame}{Méthode d'Euler}
Pour \(y'=ay+b\), avec un pas \(h>0\), on approche \(y(t_k)\) par \(Y_k\) selon \(t_{k+1}=t_k+h\), \(Y_{k+1}=Y_k+h(aY_k+b)\). C'est une approximation utilisant la tangente, pas une égalité sur la solution exacte. Un pas plus petit peut améliorer l'approximation ; contrôler numériquement l'erreur sur une solution connue.
\end{frame}
\begin{frame}{Activité de recherche}
\textbf{Refroidissement d'une boisson - 40 min.}

Une boisson initialement à \(80\,^{\circ}\mathrm C\) est placée dans une pièce maintenue à \(20\,^{\circ}\mathrm C\). Pour \(t\ge0\), on note \(T(t)\) sa température en degrés Celsius, \(t\) désignant le temps écoulé en minutes. On admet le modèle \(T' (t)=-0{,}1(T(t)-20)\) et \(T(0)=80\).
\end{frame}
\begin{frame}{Activité de recherche (2)}
\begin{enumerate}
\setcounter{enumi}{0}\item Résoudre ce problème initial et donner l'expression de \(T(t)\).

\item Étudier le sens de variation de \(T\) et déterminer sa limite en \(+\infty\). Interpréter les résultats.

\item Déterminer le premier instant où la température vaut \(40\,^{\circ}\mathrm C\). Donner la valeur exacte, puis une valeur approchée en minutes.

\item Appliquer deux étapes de la méthode d'Euler avec un pas de \(1\) minute. Comparer l'approximation obtenue au bout de deux minutes à la valeur exacte \(T(2)\).

\end{enumerate}
\end{frame}
\begin{frame}{Activité de recherche (3)}
\begin{enumerate}
\setcounter{enumi}{4}\item Indiquer une situation concrète dans laquelle l'hypothèse d'une température ambiante constante ne serait plus adaptée.

\end{enumerate}
\end{frame}
\end{document}
